\section*{Abstract}
\addcontentsline{toc}{section}{Abstract}

The southern African Cheetah (\textit{Acinonyx jubatus jubatus}), has had their range decreased drastically in the past decades. This has only been exacerbated in the past decades due to both natural and anthropogenic reasons, leaving their current extant range majority outside protected areas \parencite{weise2017distribution}. This leaves barriers like ranchers, roads and other human development in areas where cheetahs will be roaming \parencite{marker2008spatial}. Models looking at habitat suitability have been done in the past, but very little predictive modeling of the movement network permeability. One 2026 KAZA (Kavango-Zambezi Transfrontier Conservation Area) multi-species connectivity study excluded the cheetah outright \parencite{sousa2026kaza}. A similar study was completed for the Asiatic cheetah (\textit{Acinonyx jubatus venaticus}), identifying potential connectivity through the Iranian landscape \parencite{moqanaki2017iran}. This study attempts to identify core areas within southern Africa using population data from the Dryad dataset accompanying \textcite{weise2017distribution} \parencite{weise2017data}, model connectivity between them using a multivariate analysis, then see how that connectivity has changed between 2012 and 2024. The modeled connection cost changed only slightly overall. Out of the 253 core pairs, 161 became somewhat easier and 92 became harder. Outside the core areas, only about 23\% of the landscape that is most important for cheetah movement was protected. This is lower than the 31.26\% of protected land across the wider study area. Protection was higher only in the very most connected places. Of the 45 links selected by the model, 32 were kept as the main priority links. The other 13 are still possible connections, but need more field data before they can be treated as priorities.

\vspace{1em}
\noindent\textbf{Keywords:} landscape connectivity; circuit theory; \textit{Acinonyx jubatus};
protected area coverage; monitoring prioritisation; southern Africa
